% Reviewed translation: PDF pages 298--310 / printed pages 284--296.
% Reconstructed and formula-checked against the original scan.
\MathBookSourcePage{298}
\section{定常 Navier--Stokes 方程的理论}
\label{sec:navier-stokes-steady-theory}
\setcounter{equation}{0}

如 Paragraph~I.5 所述, 定常 Navier--Stokes 方程可以写成
\begin{equation}\label{eq:navier-steady-system}
\left\{
\begin{aligned}
-\nu\Delta\mathbf u
+\sum_{j=1}^{N}u_j\frac{\partial\mathbf u}{\partial x_j}
+\operatorname{grad}p&=\mathbf f&&\text{in }\Omega,\\
\operatorname{div}\mathbf u&=0&&\text{in }\Omega,
\end{aligned}
\right.
\qquad
\mathbf f\text{ given in }H^{-1}(\Omega)^N.
\end{equation}
这里仍设 $\Omega$ 是 $\mathbb R^N$ ($N=2,3$) 中具有 Lipschitz 连续边界
$\Gamma$ 的有界区域. 本节的目标是从 Section~\ref{sec:navier-stokes-nonlinear-problems}
的抽象结果推出 Navier--Stokes 方程各种表述的存在唯一性结果.

\subsection{速度--压力表述中的 Dirichlet 问题}
\label{subsec:navier-steady-velocity-pressure}

先考虑齐次 Dirichlet 边界条件
\begin{equation}\label{eq:navier-steady-homogeneous-dirichlet}
\mathbf u=0\quad\text{on }\Gamma.
\end{equation}
为把问题 \eqref{eq:navier-steady-system},
\eqref{eq:navier-steady-homogeneous-dirichlet} 写成变分形式, 引入三线性形式
\begin{equation}\label{eq:navier-steady-convection-form}
a_1(\mathbf w;\mathbf u,\mathbf v)
=\sum_{i,j=1}^{N}\int_\Omega
w_j\frac{\partial u_i}{\partial x_j}v_i\,dx.
\end{equation}
下面两个 Lemma 给出三线性形式 $a_1(\cdot;\cdot,\cdot)$ 的有用性质.

\setcounter{statement}{0}
\begin{Lemma}\label{lem:navier-steady-convection-continuity}
当 $N\leq4$ 时, 三线性形式 $a_1(\cdot;\cdot,\cdot)$ 在
$(H^1(\Omega)^N)^3$ 上连续.
\end{Lemma}
\begin{Proof}
由 Sobolev 嵌入 Theorem~\ref{thm:sobolev-embedding}, 当 $N\leq4$ 时,
$H^1(\Omega)$ 连续嵌入 $L^4(\Omega)$. 因而由 H\"older 不等式,
若 $\mathbf u,\mathbf v,\mathbf w\in H^1(\Omega)^N$, 则
\[
w_j\frac{\partial u_i}{\partial x_j}v_i\in L^1(\Omega),
\qquad 1\leq i,j\leq N,
\]
并且
\[
\begin{aligned}
\left|\int_\Omega
w_j\frac{\partial u_i}{\partial x_j}v_i\,dx\right|
&\leq
\|w_j\|_{0,4,\Omega}
\left\|\frac{\partial u_i}{\partial x_j}\right\|_{0,\Omega}
\|v_i\|_{0,4,\Omega}\\
&\leq C_1\|w_j\|_{1,\Omega}\|u_i\|_{1,\Omega}\|v_i\|_{1,\Omega}.
\end{aligned}
\]
所以 $a_1(\cdot;\cdot,\cdot)$ 定义良好且连续, 并满足
\[
|a_1(\mathbf w;\mathbf u,\mathbf v)|
\leq C_2\|\mathbf u\|_{1,\Omega}
\|\mathbf v\|_{1,\Omega}\|\mathbf w\|_{1,\Omega}.
\]
\end{Proof}

\MathBookSourcePage{299}
\setcounter{statement}{1}
\begin{Lemma}\label{lem:navier-steady-convection-skew}
设 $\mathbf u,\mathbf v\in H^1(\Omega)^N$, 且
$\mathbf w\in H^1(\Omega)^N$ 满足
$\operatorname{div}\mathbf w=0$ 和
$\mathbf w\cdot\mathbf n|_\Gamma=0$. 则
\begin{align}
a_1(\mathbf w;\mathbf u,\mathbf v)
+a_1(\mathbf w;\mathbf v,\mathbf u)&=0,
\label{eq:navier-steady-convection-antisymmetry}\\
a_1(\mathbf w;\mathbf v,\mathbf v)&=0.
\label{eq:navier-steady-convection-cancellation}
\end{align}
\end{Lemma}
\begin{Proof}
性质 \eqref{eq:navier-steady-convection-antisymmetry} 和
\eqref{eq:navier-steady-convection-cancellation} 显然等价, 因而只需验证后者.
令 $\mathbf v\in\mathcal D(\overline\Omega)^N$ 且
$\mathbf w\in H^1(\Omega)^N$. 由 Green 公式 \eqref{eq:sec2-hdiv-green},
\[
\begin{aligned}
a_1(\mathbf w;\mathbf v,\mathbf v)
&=\frac12\sum_{i,j=1}^{N}\int_\Omega
w_j\frac{\partial(v_i^2)}{\partial x_j}\,dx\\
&=-\frac12\int_\Omega(\operatorname{div}\mathbf w)|\mathbf v|^2\,dx
+\frac12\int_\Gamma(\mathbf w\cdot\mathbf n)|\mathbf v|^2\,ds=0.
\end{aligned}
\]
最后利用 $\mathcal D(\overline\Omega)^N$ 在 $H^1(\Omega)^N$ 中的稠密性
(Theorem~\ref{thm:sobolev-density-extension}) 即得结论.
\end{Proof}

置
\[
\mathcal V=\{\mathbf v\in\mathcal D(\Omega)^N;
\operatorname{div}\mathbf v=0\},
\qquad
V=\{\mathbf v\in H_0^1(\Omega)^N;
\operatorname{div}\mathbf v=0\}.
\]
再定义
\begin{equation}\label{eq:navier-steady-viscous-form}
a_0(\mathbf u,\mathbf v)
=\nu(\operatorname{grad}\mathbf u,\operatorname{grad}\mathbf v)
=2\nu\sum_{i,j=1}^{N}(D_{ij}(\mathbf u),D_{ij}(\mathbf v)),
\end{equation}
其中只要 $\mathbf u$ 或 $\mathbf v$ 属于 $V$, 第二个等式就成立. 最后定义
\begin{equation}\label{eq:navier-steady-total-form}
a(\mathbf w;\mathbf u,\mathbf v)
=a_0(\mathbf u,\mathbf v)+a_1(\mathbf w;\mathbf u,\mathbf v).
\end{equation}
于是问题 \eqref{eq:navier-steady-system},
\eqref{eq:navier-steady-homogeneous-dirichlet} 的变分形式为: 求
$(\mathbf u,p)\in V\times L_0^2(\Omega)$, 使
\begin{equation}\label{eq:navier-steady-homogeneous-variational}
a(\mathbf u;\mathbf u,\mathbf v)
-(p,\operatorname{div}\mathbf v)
=\langle\mathbf f,\mathbf v\rangle
\qquad\forall\mathbf v\in H_0^1(\Omega)^N.
\end{equation}

\setcounter{statement}{0}
\begin{Theorem}\label{thm:navier-steady-homogeneous-existence}
设 $N\leq4$, $\Omega$ 是具有 Lipschitz 连续边界的有界区域, 且
$\mathbf f\in H^{-1}(\Omega)^N$. 则问题
\eqref{eq:navier-steady-homogeneous-variational} 至少有一个解
$(\mathbf u,p)$. 此外, 该问题等价于
\eqref{eq:navier-steady-system},
\eqref{eq:navier-steady-homogeneous-dirichlet}.
\end{Theorem}

\begin{Proof}
\MathBookSourcePage{300}
在 Section~\ref{sec:navier-stokes-nonlinear-problems} 的抽象框架中取
\[
X=H_0^1(\Omega)^N,\qquad
\|\mathbf v\|_X=|\mathbf v|_{1,\Omega},\qquad
M=L_0^2(\Omega),
\]
并令
\[
b(\mathbf v,q)=-(q,\operatorname{div}\mathbf v),
\qquad
\langle l,\mathbf v\rangle=\langle\mathbf f,\mathbf v\rangle.
\]
应用 Theorem~\ref{thm:navier-abstract-existence} 和
Theorem~\ref{thm:navier-abstract-lagrange-multiplier}. 由
Lemma~\ref{lem:navier-steady-convection-skew},
\[
a(\mathbf w;\mathbf v,\mathbf v)
=a_0(\mathbf v,\mathbf v)=\nu|\mathbf v|_{1,\Omega}^2.
\]
还需验证映射 $\mathbf w\mapsto A(\mathbf w)\mathbf w$ 从 $V$ 到 $V'$ 弱连续.
事实上, 若 $\mathbf w_n\rightharpoonup\mathbf w$ 于 $V$, 则由紧嵌入
$H_0^1(\Omega)\hookrightarrow L^2(\Omega)$,
$\mathbf w_n\to\mathbf w$ 于 $L^2(\Omega)^N$. 因此对每个
$\mathbf v\in\mathcal V$, 有
\[
a_1(\mathbf w_n;\mathbf w_n,\mathbf v)
\longrightarrow a_1(\mathbf w;\mathbf w,\mathbf v).
\]
再利用 $\mathcal V$ 在 $V$ 中的稠密性
(Corollary~\ref{cor:sec2-test-divfree-dense})和
Lemma~\ref{lem:navier-steady-convection-continuity} 的一致界, 即得所需弱连续性.

所以存在 $\mathbf u\in V$ 满足
\begin{equation}\label{eq:navier-steady-reduced-variational}
a(\mathbf u;\mathbf u,\mathbf v)
=\langle\mathbf f,\mathbf v\rangle
\qquad\forall\mathbf v\in V.
\end{equation}
由压力 inf--sup 不等式 \eqref{eq:stokes-pressure-infsup}, 存在唯一的
$p\in L_0^2(\Omega)$ 使 $(\mathbf u,p)$ 满足
\eqref{eq:navier-steady-homogeneous-variational}.
该变分问题与原边值问题的等价性由分布意义下的方程和迹条件直接得到.
\end{Proof}

为陈述唯一性条件, 置
\begin{equation}\label{eq:navier-steady-convection-constant}
\mathcal N
=\sup_{\mathbf u,\mathbf v,\mathbf w\in V}
\frac{a_1(\mathbf w;\mathbf u,\mathbf v)}
{|\mathbf u|_{1,\Omega}|\mathbf v|_{1,\Omega}|\mathbf w|_{1,\Omega}},
\end{equation}
以及
\begin{equation}\label{eq:navier-steady-force-dual-norm}
\|\mathbf f\|_{V'}
=\sup_{\mathbf v\in V}
\frac{\langle\mathbf f,\mathbf v\rangle}{|\mathbf v|_{1,\Omega}}.
\end{equation}

\MathBookSourcePage{301}
\setcounter{statement}{1}
\begin{Theorem}\label{thm:navier-steady-homogeneous-uniqueness}
在 Theorem~\ref{thm:navier-steady-homogeneous-existence} 的假设下, 若
\begin{equation}\label{eq:navier-steady-small-data-condition}
\frac{\mathcal N}{\nu^2}\|\mathbf f\|_{V'}<1,
\end{equation}
则问题 \eqref{eq:navier-steady-homogeneous-variational} 的解唯一.
\end{Theorem}
\begin{Proof}
应用 Theorem~\ref{thm:navier-abstract-uniqueness}.
由 \eqref{eq:navier-steady-convection-constant}, 对任意
$\mathbf u,\mathbf v,\mathbf w\in V$ 有
\[
|a(\mathbf w;\mathbf u,\mathbf v)-a(\mathbf u;\mathbf u,\mathbf v)|
=|a_1(\mathbf w-\mathbf u;\mathbf u,\mathbf v)|
\leq\mathcal N|\mathbf w-\mathbf u|_{1,\Omega}
|\mathbf u|_{1,\Omega}|\mathbf v|_{1,\Omega}.
\]
抽象唯一性条件恰为 \eqref{eq:navier-steady-small-data-condition}.
\end{Proof}

现在考虑非齐次 Dirichlet 边界条件
\begin{equation}\label{eq:navier-steady-nonhomogeneous-dirichlet}
\mathbf u=\mathbf g\quad\text{on }\Gamma.
\end{equation}
设 $\Gamma_i$, $0\leq i\leq p$, 是 $\Gamma$ 的连通分支. 假定
\begin{equation}\label{eq:navier-steady-component-flux}
\int_{\Gamma_i}\mathbf g\cdot\mathbf n\,ds=0,
\qquad 0\leq i\leq p.
\end{equation}
下面给出 Hopf~\cite{Hopf1955} 的一个基本结果.

\setcounter{statement}{2}
\begin{Lemma}\label{lem:navier-steady-hopf-lifting}
设 $N\leq3$, $\mathbf g\in H^{1/2}(\Gamma)^N$ 且满足
\eqref{eq:navier-steady-component-flux}. 对每个 $\varepsilon>0$, 存在
$\mathbf u_0(\varepsilon)\in H^1(\Omega)^N$, 使
\begin{equation}\label{eq:navier-steady-hopf-lifting-properties}
\operatorname{div}\mathbf u_0=0\quad\text{in }\Omega,
\qquad
\mathbf u_0|_\Gamma=\mathbf g,
\end{equation}
并且
\begin{equation}\label{eq:navier-steady-hopf-small-convection}
|a_1(\mathbf v;\mathbf u_0,\mathbf v)|
\leq\varepsilon|\mathbf v|_{1,\Omega}^2
\qquad\forall\mathbf v\in V.
\end{equation}
\end{Lemma}

其证明需要下面两个 Lemma.

\MathBookSourcePage{302}
\setcounter{statement}{3}
\begin{Lemma}\label{lem:navier-steady-boundary-cutoff}
对每个 $\varepsilon>0$, 存在 $\theta_\varepsilon\in C^2(\overline\Omega)$, 使
\begin{equation}\label{eq:navier-steady-boundary-cutoff-properties}
\left\{
\begin{aligned}
\theta_\varepsilon(x)&=1&&\text{在 }\Gamma\text{ 的一个邻域中},\\
\theta_\varepsilon(x)&=0&&\text{若 }d(x;\Gamma)\geq2\delta(\varepsilon),\\
\left|\frac{\partial\theta_\varepsilon}{\partial x_i}(x)\right|
&\leq\frac{\varepsilon}{d(x;\Gamma)}
&&\text{若 }d(x;\Gamma)\leq2\delta(\varepsilon),\quad1\leq i\leq N,
\end{aligned}
\right.
\qquad
\delta(\varepsilon)=\exp(-1/\varepsilon).
\end{equation}
\end{Lemma}
\begin{Proof}
先定义
\[
\phi_\varepsilon(\mu)=
\begin{cases}
1,&0\leq\mu\leq\delta(\varepsilon)^2,\\
\varepsilon\log\bigl(\delta(\varepsilon)/\mu\bigr),
&\delta(\varepsilon)^2\leq\mu\leq\delta(\varepsilon),\\
0,&\mu\geq\delta(\varepsilon),
\end{cases}
\]
并令 $\chi_\varepsilon(x)=\phi_\varepsilon(d(x;\Gamma))$.
对 $\chi_\varepsilon$ 作适当正则化, 即得到满足
\eqref{eq:navier-steady-boundary-cutoff-properties} 的函数
$\theta_\varepsilon$.
\end{Proof}

\setcounter{statement}{4}
\begin{Lemma}\label{lem:navier-steady-hardy-boundary}
存在常数 $C$, 使
\begin{equation}\label{eq:navier-steady-hardy-boundary}
\left\|\frac{\phi}{d(\cdot;\Gamma)}\right\|_{0,\Omega}
\leq C|\phi|_{1,\Omega}
\qquad\forall\phi\in H_0^1(\Omega).
\end{equation}
\end{Lemma}
\begin{Proof}
借助有限个局部坐标图把 $\Omega$ 的边界邻域化为半空间, 便可把结论化为一维不等式
\begin{equation}\label{eq:navier-steady-hardy-one-dimensional}
\int_0^\infty\left|\frac{\phi(t)}{t}\right|^2dt
\leq4\int_0^\infty|\phi'(t)|^2dt
\qquad\forall\phi\in\mathcal D(0,\infty).
\end{equation}
\MathBookSourcePage{303}
为证明 \eqref{eq:navier-steady-hardy-one-dimensional}, 先写成
\[
\phi(t)=\int_0^t\phi'(s)\,ds,
\]
再置 $t=e^\tau$, $s=e^\sigma$. 于是相关积分可写成含经典
Heaviside 函数 $H$ 的卷积. 应用标准卷积不等式
\[
\|f*g\|_2\leq\|f\|_1\|g\|_2,
\]
得到
\[
\int_{-\infty}^{\infty}e^{-\tau}|\phi(e^\tau)|^2\,d\tau
\leq4\int_{-\infty}^{\infty}e^\tau|\phi'(e^\tau)|^2\,d\tau.
\]
变回变量 $t$ 即为 \eqref{eq:navier-steady-hardy-one-dimensional}, 从而得到结论.
\end{Proof}

现在证明 Lemma~\ref{lem:navier-steady-hopf-lifting}.
\begin{Proof}
由 Lemma~\ref{lem:sec2-div-free-lifting}, 存在
$\mathbf w_0\in H^1(\Omega)^N$, 使
$\operatorname{div}\mathbf w_0=0$ 且 $\mathbf w_0|_\Gamma=\mathbf g$.
条件 \eqref{eq:navier-steady-component-flux}、
Theorem~\ref{thm:sec3-stream-function-characterization} 及其三维对应结果给出
$\boldsymbol\psi_0\in H^2(\Omega)^N$, 使
\[
\mathbf w_0=\operatorname{curl}\boldsymbol\psi_0.
\]
二维情形把 $\boldsymbol\psi_0$ 理解为标量流函数. 对 $N=3$, 定义
\[
\mathbf u_{0\mu}
=\operatorname{curl}(\theta_\mu\boldsymbol\psi_0).
\]
则
\[
\operatorname{div}\mathbf u_{0\mu}=0,
\qquad
\mathbf u_{0\mu}|_\Gamma=\mathbf g.
\]
二维情形完全类似.

\MathBookSourcePage{304}
由 Lemma~\ref{lem:navier-steady-boundary-cutoff},
\[
|\mathbf u_{0\mu}(x)|
\leq C_1\left\{
\frac{\mu}{d(x;\Gamma)}|\boldsymbol\psi_0(x)|
+|D\boldsymbol\psi_0(x)|\right\},
\]
其中
\[
|D\boldsymbol\psi_0|^2
=\sum_{i,j=1}^{N}
\left|\frac{\partial\psi_{0i}}{\partial x_j}\right|^2.
\]
由于 $H^2(\Omega)\hookrightarrow C^0(\overline\Omega)$, 再应用
Lemma~\ref{lem:navier-steady-hardy-boundary}、H\"older 不等式和
$H^1(\Omega)\hookrightarrow L^6(\Omega)$, 对
$\mathbf v\in V$ 得
\[
\|v_i u_{0\mu,j}\|_{0,\Omega}
\leq C_5\{\mu+\varphi(\mu)\}|v_i|_{1,\Omega},
\]
其中
\[
\varphi(\mu)=
\left(\int_{\{d(x;\Gamma)\leq2\delta(\mu)\}}
|D\boldsymbol\psi_0(x)|^3\,dx\right)^{1/3}.
\]
由 Lemma~\ref{lem:navier-steady-convection-skew},
\[
|a_1(\mathbf v;\mathbf u_{0\mu},\mathbf v)|
=|a_1(\mathbf v;\mathbf v,\mathbf u_{0\mu})|
\leq C_6\{\mu+\varphi(\mu)\}|\mathbf v|_{1,\Omega}^2.
\]
当 $\mu\to0$ 时 $\varphi(\mu)\to0$. 选择 $\mu$ 使
$C_6\{\mu+\varphi(\mu)\}\leq\varepsilon$, 相应的
$\mathbf u_0=\mathbf u_{0\mu}$ 即满足要求.
\end{Proof}

\MathBookSourcePage{305}
非齐次问题的变分表述为: 求
$(\mathbf u,p)\in H^1(\Omega)^N\times L_0^2(\Omega)$, 使
\begin{equation}\label{eq:navier-steady-nonhomogeneous-variational}
\left\{
\begin{aligned}
a(\mathbf u;\mathbf u,\mathbf v)
-(p,\operatorname{div}\mathbf v)
&=\langle\mathbf f,\mathbf v\rangle
&&\forall\mathbf v\in H_0^1(\Omega)^N,\\
\operatorname{div}\mathbf u&=0&&\text{in }\Omega,\\
\mathbf u&=\mathbf g&&\text{on }\Gamma.
\end{aligned}
\right.
\end{equation}

\setcounter{statement}{2}
\begin{Theorem}\label{thm:navier-steady-nonhomogeneous-existence}
设 $N\leq3$, $\Omega$ 是具有 Lipschitz 连续边界的有界区域,
$\mathbf f\in H^{-1}(\Omega)^N$, 且
$\mathbf g\in H^{1/2}(\Gamma)^N$ 满足
\eqref{eq:navier-steady-component-flux}. 则问题
\eqref{eq:navier-steady-nonhomogeneous-variational} 至少有一个解.
\end{Theorem}
\begin{Proof}
由 Lemma~\ref{lem:navier-steady-hopf-lifting} 选取
$\mathbf u_0$ 并写成 $\mathbf u=\mathbf u_0+\mathbf w$, 其中
$\mathbf w\in V$. 展开得
\[
\begin{aligned}
a(\mathbf u_0+\mathbf w;\mathbf u_0+\mathbf w,\mathbf v)
={}&a(\mathbf w;\mathbf w,\mathbf v)
+a_1(\mathbf u_0;\mathbf w,\mathbf v)
+a_1(\mathbf w;\mathbf u_0,\mathbf v)\\
&+a(\mathbf u_0;\mathbf u_0,\mathbf v).
\end{aligned}
\]
定义
\[
\widetilde a(\mathbf w;\mathbf u,\mathbf v)
=a(\mathbf w;\mathbf u,\mathbf v)
+a_1(\mathbf u_0;\mathbf u,\mathbf v)
+a_1(\mathbf u;\mathbf u_0,\mathbf v).
\]
于是只需解
\begin{equation}\label{eq:navier-steady-lifted-variational}
\widetilde a(\mathbf w;\mathbf w,\mathbf v)
-(p,\operatorname{div}\mathbf v)
=\langle\mathbf f,\mathbf v\rangle
-a(\mathbf u_0;\mathbf u_0,\mathbf v)
\quad\forall\mathbf v\in H_0^1(\Omega)^N.
\end{equation}
取 $\varepsilon<\nu$. 由
\eqref{eq:navier-steady-hopf-small-convection} 和
Lemma~\ref{lem:navier-steady-convection-skew},
\[
\widetilde a(\mathbf w;\mathbf v,\mathbf v)
=\nu|\mathbf v|_{1,\Omega}^2
+a_1(\mathbf v;\mathbf u_0,\mathbf v)
\geq(\nu-\varepsilon)|\mathbf v|_{1,\Omega}^2.
\]
\MathBookSourcePage{306}
因此 Theorem~\ref{thm:navier-abstract-existence} 适用于
\eqref{eq:navier-steady-lifted-variational}, 给出至少一个解.
\end{Proof}

为得到非齐次问题的唯一性条件, 对满足
\eqref{eq:navier-steady-hopf-lifting-properties} 的 $\mathbf u_0$ 定义
\begin{equation}\label{eq:navier-steady-lifting-rho}
\rho(\mathbf u_0)
=\sup_{\mathbf v\in V}
\frac{a_1(\mathbf v;\mathbf u_0,\mathbf v)}{|\mathbf v|_{1,\Omega}^2},
\end{equation}
以及
\begin{equation}\label{eq:navier-steady-lifted-load-norm}
\|l(\mathbf f;\mathbf u_0)\|_{V'}
=\sup_{\mathbf v\in V}
\frac{\langle l(\mathbf f;\mathbf u_0),\mathbf v\rangle}
{|\mathbf v|_{1,\Omega}},
\end{equation}
其中
\[
\langle l(\mathbf f;\mathbf u_0),\mathbf v\rangle
=\langle\mathbf f,\mathbf v\rangle
-a(\mathbf u_0;\mathbf u_0,\mathbf v).
\]
最后置
\begin{equation}\label{eq:navier-steady-critical-viscosity}
\nu_0=
\inf\left\{
\rho(\mathbf u_0)
+\bigl(\mathcal N\|l(\mathbf f;\mathbf u_0)\|_{V'}\bigr)^{1/2};
\ \mathbf u_0\in H^1(\Omega)^N
\text{ satisfies \eqref{eq:navier-steady-hopf-lifting-properties}}
\right\}.
\end{equation}

\setcounter{statement}{3}
\begin{Theorem}\label{thm:navier-steady-nonhomogeneous-uniqueness}
在 Theorem~\ref{thm:navier-steady-nonhomogeneous-existence} 的假设下,
若 $\nu>\nu_0$, 则问题
\eqref{eq:navier-steady-nonhomogeneous-variational} 的解唯一.
\end{Theorem}
\begin{Proof}
由 \eqref{eq:navier-steady-critical-viscosity}, 可选择一个满足
\eqref{eq:navier-steady-hopf-lifting-properties} 的 $\mathbf u_0$, 使
\[
\rho(\mathbf u_0)
+\bigl(\mathcal N\|l(\mathbf f;\mathbf u_0)\|_{V'}\bigr)^{1/2}<\nu.
\]
将 Theorem~\ref{thm:navier-abstract-uniqueness} 应用于
\eqref{eq:navier-steady-lifted-variational} 即得结论.
\end{Proof}

\MathBookSourcePage{307}
\setcounter{statement}{0}
\begin{Remark}\label{rem:navier-steady-lifted-bound}
上述证明还给出
\[
|\mathbf w|_{1,\Omega}
\leq\frac{\|l(\mathbf f;\mathbf u_0)\|_{V'}}
{\nu-\rho(\mathbf u_0)}.
\]
\end{Remark}

\setcounter{statement}{1}
\begin{Remark}\label{rem:navier-steady-oseen-iteration}
可以用 Section~\ref{sec:navier-stokes-nonlinear-problems} 的逐次逼近构造解.
在速度--压力变量中, 该迭代等价于
\begin{equation}\label{eq:navier-steady-oseen-iteration}
\left\{
\begin{aligned}
-\nu\Delta\mathbf u^{m+1}
+\sum_{j=1}^{N}u_j^m
\frac{\partial\mathbf u^{m+1}}{\partial x_j}
+\operatorname{grad}p^{m+1}
&=\mathbf f&&\text{in }\Omega,\\
\operatorname{div}\mathbf u^{m+1}&=0&&\text{in }\Omega,\\
\mathbf u^{m+1}&=\mathbf g&&\text{on }\Gamma.
\end{aligned}
\right.
\end{equation}
在 Theorem~\ref{thm:navier-steady-nonhomogeneous-uniqueness} 的条件下,
该序列收敛到唯一解.
\end{Remark}

\subsection{齐次 Dirichlet 问题的流函数表述}
\label{subsec:navier-steady-stream-function}

先设 $N=2$. 由 Chapter~\ref{chap:stokes-foundations} 的结果, 每个
$\mathbf u\in V$ 都可唯一写成 $\mathbf u=\operatorname{curl}\psi$, 其中
\[
\Psi=\left\{\chi\in H^2(\Omega);
\ \chi|_{\Gamma_0}=0,
\ \chi|_{\Gamma_i}\text{ is constant},\ 1\leq i\leq p,
\ \frac{\partial\chi}{\partial n}\bigg|_\Gamma=0\right\}.
\]
直接计算给出
\begin{equation}\label{eq:navier-steady-two-dimensional-curl-convection}
a_1(\mathbf u;\mathbf u,\mathbf v)
=\int_\Omega\operatorname{curl}\mathbf u\,(u_1v_2-u_2v_1)\,dx.
\end{equation}
其中 $\mathbf u\in H_0^1(\Omega)^2$ 且 $\mathbf v\in V$.
事实上,
\begin{equation}\label{eq:navier-steady-two-dimensional-component-identities}
\left\{
\begin{aligned}
u_1\frac{\partial u_1}{\partial x_1}
+u_2\frac{\partial u_1}{\partial x_2}
&=\frac{\partial}{\partial x_1}\frac{|\mathbf u|^2}{2}
-u_2\operatorname{curl}\mathbf u,\\
u_1\frac{\partial u_2}{\partial x_1}
+u_2\frac{\partial u_2}{\partial x_2}
&=\frac{\partial}{\partial x_2}\frac{|\mathbf u|^2}{2}
+u_1\operatorname{curl}\mathbf u.
\end{aligned}
\right.
\end{equation}
再分部积分消去 $\operatorname{grad}(|\mathbf u|^2)$ 项, 即得到
\eqref{eq:navier-steady-two-dimensional-curl-convection}.
\MathBookSourcePage{308}
若 $\mathbf u=\operatorname{curl}\psi$ 且
$\mathbf v=\operatorname{curl}\phi$, 则
\begin{equation}\label{eq:navier-steady-two-dimensional-stream-convection}
a_1(\mathbf u;\mathbf u,\mathbf v)
=\int_\Omega\Delta\psi
\left(
\frac{\partial\psi}{\partial x_2}
\frac{\partial\phi}{\partial x_1}
-\frac{\partial\psi}{\partial x_1}
\frac{\partial\phi}{\partial x_2}
\right)dx.
\end{equation}
又由 Theorem~\ref{thm:stokes-two-dimensional-stream-function} 证明中的恒等式
\eqref{eq:stokes-stream-gradient-laplacian-identity},
\begin{equation}\label{eq:navier-steady-two-dimensional-stream-form}
\begin{aligned}
a(\mathbf u;\mathbf u,\mathbf v)
={}&\nu(\Delta\psi,\Delta\phi)\\
&+\int_\Omega\Delta\psi
\left(
\frac{\partial\psi}{\partial x_2}
\frac{\partial\phi}{\partial x_1}
-\frac{\partial\psi}{\partial x_1}
\frac{\partial\phi}{\partial x_2}
\right)dx.
\end{aligned}
\end{equation}
因此二维问题等价于如下非线性双调和问题: 求 $\psi\in\Psi$, 使
\[
\nu(\Delta\psi,\Delta\phi)
+\int_\Omega\Delta\psi
\left(
\frac{\partial\psi}{\partial x_2}
\frac{\partial\phi}{\partial x_1}
-\frac{\partial\psi}{\partial x_1}
\frac{\partial\phi}{\partial x_2}
\right)dx
=\langle\mathbf f,\operatorname{curl}\phi\rangle
\quad\forall\phi\in\Psi.
\]
由于这一等价性, 前面关于存在性和唯一性的全部结果都适用于该问题.
相应强形式为
\[
\nu\Delta^2\psi
-\frac{\partial}{\partial x_1}
\left(\Delta\psi\frac{\partial\psi}{\partial x_2}\right)
+\frac{\partial}{\partial x_2}
\left(\Delta\psi\frac{\partial\psi}{\partial x_1}\right)
=\operatorname{curl}\mathbf f
\quad\text{in }\Omega,
\]
并带有边界条件
\[
\psi=0\text{ on }\Gamma_0,
\quad
\psi=\text{constant on }\Gamma_i\ (1\leq i\leq p),
\quad
\frac{\partial\psi}{\partial n}\bigg|_\Gamma=0,
\]
以及
\[
\int_{\Gamma_i}
\left(\nu\frac{\partial(\Delta\psi)}{\partial n}
-\mathbf f\cdot\boldsymbol\tau\right)ds=0,
\qquad1\leq i\leq p.
\]

现在设 $N=3$, 且 $\Omega$ 单连通. 引入空间
\MathBookSourcePage{309}
\[
\begin{aligned}
\Psi=\{\boldsymbol\phi\in L^2(\Omega)^3;\ &
\operatorname{div}\boldsymbol\phi\in H^1(\Omega),
\ \operatorname{curl}\boldsymbol\phi\in H_0^1(\Omega)^3,\\
&\boldsymbol\phi\times\mathbf n|_\Gamma=0,
\ \int_{\Gamma_i}\boldsymbol\phi\cdot\mathbf n\,ds=0,
\ 0\leq i\leq p\}.
\end{aligned}
\]
其范数为
\[
\|\boldsymbol\phi\|
=\left\{\|\boldsymbol\phi\|_{0,\Omega}^2
+\|\operatorname{div}\boldsymbol\phi\|_{1,\Omega}^2
+\|\operatorname{curl}\boldsymbol\phi\|_{1,\Omega}^2\right\}^{1/2}.
\]
由 Lemma~\ref{lem:stokes-vector-biharmonic-norm-tangential}, 该范数等价于
$\|\Delta\boldsymbol\phi\|_{0,\Omega}$; 以下把后者简记为
$|\boldsymbol\phi|$.

每个 $\mathbf u\in V$ 都有唯一的向量势
$\boldsymbol\psi\in\Psi$, 使
$\mathbf u=\operatorname{curl}\boldsymbol\psi$ 且
$\operatorname{div}\boldsymbol\psi=0$. 恒等式
\begin{equation}\label{eq:navier-steady-three-dimensional-convection-identity}
\sum_{i,j=1}^{3}u_j\frac{\partial u_i}{\partial x_j}v_i
=(\operatorname{curl}\mathbf u\times\mathbf u)\cdot\mathbf v
+\frac12\operatorname{grad}(|\mathbf u|^2)\cdot\mathbf v,
\end{equation}
给出流函数问题: 求 $\boldsymbol\psi\in\Psi$, 使
\begin{equation}\label{eq:navier-steady-three-dimensional-stream-problem}
\nu(\Delta\boldsymbol\psi,\Delta\boldsymbol\phi)
-\int_\Omega
(\Delta\boldsymbol\psi\times\operatorname{curl}\boldsymbol\psi)
\cdot\operatorname{curl}\boldsymbol\phi\,dx
=\langle\mathbf f,\operatorname{curl}\boldsymbol\phi\rangle
\quad\forall\boldsymbol\phi\in\Psi.
\end{equation}
由 Theorem~\ref{thm:navier-steady-homogeneous-existence}, 问题
\eqref{eq:navier-steady-three-dimensional-stream-problem} 至少有一个满足
$\operatorname{div}\boldsymbol\psi=0$ 的解. 反过来, 对任一解置
$\mathbf w=\operatorname{curl}\boldsymbol\psi$. 由原书所引用的
Problem \eqref{eq:navier-steady-system}, \eqref{eq:navier-steady-convection-form}
的计算可得
\[
\begin{aligned}
\nu(\operatorname{grad}\mathbf w,\operatorname{grad}\mathbf v)
+\int_\Omega\operatorname{curl}\mathbf w\times\mathbf w\cdot\mathbf v\,dx
={}&\langle\mathbf f,\mathbf v\rangle\\
&+\int_\Omega
\operatorname{grad}(\operatorname{div}\boldsymbol\psi)
\times\mathbf w\cdot\mathbf v\,dx.
\end{aligned}
\]
除非最后一个积分为零(而该性质并不能由
\eqref{eq:navier-steady-three-dimensional-stream-problem} 推出),
$\mathbf w$ 不满足 Navier--Stokes 方程.

定义
\begin{equation}\label{eq:navier-steady-three-dimensional-trilinear-form}
a_1(\boldsymbol\phi;\boldsymbol\psi,\boldsymbol\chi)
=-\int_\Omega
(\Delta\boldsymbol\phi\times\operatorname{curl}\boldsymbol\psi)
\cdot\operatorname{curl}\boldsymbol\chi\,dx.
\end{equation}
该形式显然满足 \eqref{eq:navier-steady-convection-antisymmetry} 和
\eqref{eq:navier-steady-convection-cancellation} 的对应性质.
\MathBookSourcePage{310}
由 H\"older 不等式和 Sobolev 嵌入,
\[
\begin{aligned}
|a_1(\boldsymbol\phi;\boldsymbol\psi,\boldsymbol\chi)|
&\leq\|\Delta\boldsymbol\phi\|_{0,\Omega}
\|\operatorname{curl}\boldsymbol\psi\|_{0,4,\Omega}
\|\operatorname{curl}\boldsymbol\chi\|_{0,4,\Omega}\\
&\leq C_1\|\boldsymbol\phi\|
\|\boldsymbol\psi\|\|\boldsymbol\chi\|.
\end{aligned}
\]
即
\begin{equation}\label{eq:navier-steady-three-dimensional-trilinear-bound}
|a_1(\boldsymbol\phi;\boldsymbol\psi,\boldsymbol\chi)|
\leq C_1|\boldsymbol\phi||\boldsymbol\psi||\boldsymbol\chi|.
\end{equation}
再令
\[
a(\boldsymbol\phi;\boldsymbol\psi,\boldsymbol\chi)
=\nu(\Delta\boldsymbol\psi,\Delta\boldsymbol\chi)
+a_1(\boldsymbol\phi;\boldsymbol\psi,\boldsymbol\chi)
\]
等价关系
\[
\|\boldsymbol\phi\|\simeq
\|\Delta\boldsymbol\phi\|_{0,\Omega}=|\boldsymbol\phi|
\qquad\forall\boldsymbol\phi\in\Psi
\]
保证该形式的椭圆性.
以及
\begin{equation}\label{eq:navier-steady-three-dimensional-convection-constant}
\widetilde{\mathcal N}
=\sup_{\boldsymbol\phi,\boldsymbol\psi,\boldsymbol\chi\in\Psi}
\frac{a_1(\boldsymbol\phi;\boldsymbol\psi,\boldsymbol\chi)}
{|\boldsymbol\phi||\boldsymbol\psi||\boldsymbol\chi|}.
\end{equation}
该形式满足 Section~\ref{sec:navier-stokes-nonlinear-problems} 的连续性与
Lipschitz 条件. 事实上,
\[
\begin{aligned}
&|a_1(\boldsymbol\phi_1;\boldsymbol\psi,\boldsymbol\chi)
-a_1(\boldsymbol\phi_2;\boldsymbol\psi,\boldsymbol\chi)|\\
&\qquad=|a_1(\boldsymbol\phi_1-\boldsymbol\phi_2;
\boldsymbol\psi,\boldsymbol\chi)|
\leq\widetilde{\mathcal N}
|\boldsymbol\phi_1-\boldsymbol\phi_2|
|\boldsymbol\psi||\boldsymbol\chi|.
\end{aligned}
\]
所以 $a_1$ 满足 \eqref{eq:navier-abstract-local-lipschitz}, 且
$L(\mu)=\widetilde{\mathcal N}$, 与 $\mu$ 无关. 因而当 $\nu$ 足够大或
$\|\mathbf f\|_{V'}$ 足够小时, 问题
\eqref{eq:navier-steady-three-dimensional-stream-problem} 的解唯一.
由于问题 \eqref{eq:navier-steady-system},
\eqref{eq:navier-steady-homogeneous-dirichlet} 的每个解的向量势都是
\eqref{eq:navier-steady-three-dimensional-stream-problem} 的解, 这也说明前一问题有唯一速度
$\mathbf u=\operatorname{curl}\boldsymbol\psi$. 因此必有
$\operatorname{div}\boldsymbol\psi=0$. 下面的 Theorem 汇总这些结果.

\setcounter{statement}{4}
\begin{Theorem}\label{thm:navier-steady-three-dimensional-stream-wellposed}
设 $\Omega$ 是 $\mathbb R^3$ 中具有 Lipschitz 连续边界的有界单连通区域.
则问题 \eqref{eq:navier-steady-three-dimensional-stream-problem} 至少有一个解
$\boldsymbol\psi$. 若
\begin{equation}\label{eq:navier-steady-three-dimensional-small-data}
\frac{\widetilde{\mathcal N}}{\nu^2}
\sup_{\boldsymbol\phi\in\Psi}
\frac{\langle\mathbf f,\operatorname{curl}\boldsymbol\phi\rangle}
{|\boldsymbol\phi|}<1,
\end{equation}
则该解唯一. 此外, 存在 $p\in L_0^2(\Omega)$, 使
$(\mathbf u=\operatorname{curl}\boldsymbol\psi,p)$ 是问题
\eqref{eq:navier-steady-system},
\eqref{eq:navier-steady-homogeneous-dirichlet} 的唯一解, 并且
$\operatorname{div}\boldsymbol\psi=0$ 于 $\Omega$.
\end{Theorem}

\MathBookSourcePage{311}
\setcounter{statement}{2}
\begin{Remark}\label{rem:navier-steady-normal-potential-space}
在空间
\[
\Psi_1=\left\{\boldsymbol\phi\in L^2(\Omega)^3;
\ \operatorname{div}\boldsymbol\phi\in H^1(\Omega),
\ \operatorname{curl}\boldsymbol\phi\in H_0^1(\Omega)^3,
\ \boldsymbol\phi\cdot\mathbf n|_\Gamma=0\right\}
\]
中取向量势, 可以作同样的分析并得到类似的结论.
\end{Remark}
